\section{Heyting-valued Sets}

\def\h{\mathcal{H}}

\subsection{Preliminaries}

\begin{definition}
  A \emph{Heyting-valued set \(A\)} is a carrier set \(\carr A\) equipped with a
  ``Heyting-valued relation'' \(\eq[A]{-}{-} : A×A → \h\) satisfying
  \begin{gather}
    \label{hvs:symm}  \eq x y ≤ \eq y x\\
    \label{hvs:trans} \eq x y ∧ \eq y z ≤ \eq x z
  \end{gather}

  Define \(\e x ≔ \eq x x\).
\end{definition}

\begin{definition}
  A \emph{Heyting-valued morphism \(f : A ↬ B\)} is a ``Heyting-valued
  functional relation'', meaning it is a map
  \(\carr f : \carr A × \carr B → \h\)
  satisfying
  \begin{gather}
    \label{hvm:str}    f(x,y) ≤ \e x ∧ \e y\\
    \label{hvm:well}   \eq {x'} x ∧ f(x, y) ∧ \eq y {y'} ≤ f(x', y')\\
    \label{hvm:unique} f(x, y) ∧ f(x, y') ≤ \eq y {y'}\\
    \label{hvm:total}  \e x ≤ ⋁_{y ∈ B}f(x, y)
  \end{gather}
\end{definition}

\begin{fact}
  The following hold.

  \begin{enumerate}
  \item \(f(x,y)∧\eq x {x'} = f(x',y)∧\eq x {x'}\)
  \item \(\e x = ⋁_{y∈A}\eq x y\)
  \end{enumerate}
\end{fact}

\begin{construction}
  The product \(A×B\) is given by \(\carr A × \carr B\) with equality
  \[ \eq{\p{x,y}}{\p{x',y'}} ≔ \eq x {x'} ∧ \eq y {y'}\text. \]

  The projection \(π : A×B ↬ A\) is defined by
  \[ \p{\p{x,y}, x'} ↦ \eq x {x'} ∧ \e y\text. \]
  The other projection is defined similarly
\end{construction}
\begin{proof}
  Symmetry and transitivity of the equality follow clearly.

  Likewise, the projection is clearly strict and well-defined.
  Uniqueness holds via
  \begin{align*}
    π(\p{x,y},x') ∧ π(\p{x,y},x'')
    &= \eq x {x'} ∧ \e y ∧ \eq x {x''} ∧ \e y\\
    &= \eq x {x'} ∧ \eq x {x''} ∧ \e y\\
    &≤ \eq {x'} {x''} ∧ \e y\\
    &≤ \eq {x'} {x''}\text.
  \end{align*}
  Totality holds via
  \begin{align*}
    \e{\p{x,y}} = \e x ∧ \e y = ⋁_{x'∈A}\eq x {x'} ∧ \e y\text.
  \end{align*}
\end{proof}

In the following, I will not write out these proofs. A standard reference is
TODO BorceuxV3.

\begin{construction}
  The \emph{discrete Heyting-valued set \(\c A\) on \(A\)} has \(A\) as it's
  carrier and equality defined by \(\p{x,y} ↦ ⋁\set{⊤}{x=y}\).
\end{construction}

\begin{fact}
  The empty type \(∅\) is given by the empty carrier.

  The unit type \(\mb 1\) is given by \(\{\ast\}\) with \(\e\ast = ⊤\). This is
  the discrete Heyting-valued set on \(\{\ast\}\).

  The natural numbers object is the discrete Heyting-valued set on \(ℕ\).

  The rational numbers object is the discrete Heyting-valued set on \(ℚ\).
\end{fact}
\begin{note}
  The real numbers object is not necessarily discrete.
\end{note}

\begin{construction}
  The \emph{sum \(A+B\)} is given by \(\carr A+\carr B\) with equality
  \begin{align*}
    &\eq {\inl x} {\inl x'} ≔ \eq[A] x {x'}\\
    &\eq {\inl x} {\inr y'} ≔ ⊥\\
    &\eq {\inr y} {\inr x'} ≔ ⊥\\
    &\eq {\inr y} {\inr y'} ≔ \eq[B] y {y'}
  \end{align*}

  The inclusion \(ι : A ↬ A+B\) is given by \(\p{x,t} ↦ \eq {\inl x} t\).
  The other inclusion is defined similarly.
\end{construction}

\begin{construction}
  The \emph{subobject classifier \(Ω\)} is given by \(\h\) with \(∧\) as equality.
\end{construction}

\begin{construction}
  The \emph{exponential object \(B^A\)} is given by
  \[ \set{f : \carr A × \carr B → \h}{\ref{hvm:str}-\ref{hvm:unique}} \]
  with equality
  \[ \p{f,g} ↦ ⋀_{x∈A}⋁_{y∈B}f(x,y) ∧ g(x,y)\text. \]
\end{construction}
\begin{internal}
  The equality in the internal language is
  \(\for{x∈A}{\exist{y∈B}{f(x, y) ∧ g(x, y)}}\) which can be rephrased as
  \(\for{x∈A}{f(a) = g(a)}\). Thus equality models function extensionality,
  which is know to hold in Grothendieck toposes.
\end{internal}

\begin{fact}
  The \emph{power-object \(𝒫A = Ω^A\)} can be given as the discrete Heyting-valued set
  on \(\set{P : \carr A → Ω}{\eq {x'} x ∧ P(x) ≤ P(x')}\).
\end{fact}

% \begin{construction}
%   The \emph{restriction \(A\res p\)} is given by \(\set{x ∈ A}{\e x \between p}\)
% \end{construction}

\begin{construction}
  Given a Heyting-valued set \(A\) and a relation \(R : \carr A × \carr A → \h\)
  satisfying
  \begin{gather}
    \e x ≤ R x x\\
    R x y ≤ R y x\\
    R x y ∧ R y z ≤ R x z
  \end{gather}
  we can form the \emph{quotient \(\quot A R\)} with carrier
  \(\carr A\) and equality \(R\).
\end{construction}
Subquotients can be formed similarly, by merely dropping the first condition on
the relation.


\subsubsection{Internal logic}

The internal logic is as in~\cite{Sco79}. In short it is intuitionistic
higher-order logic with partial elements and equality. Furthermore, every
Heyting-valued relation \(R : \carr A × \carr B → \h\) (satisfying
\ref{hvm:str}--\ref{hvm:well}) is a relational symbol in the internal language.

\begin{definition}
  The \emph{interpretation \(\i φ\)} of a formula \(φ\) is defined recursively
  on the structure of the formula.

  TODO: ja
\end{definition}

\begin{example}
  We can work out an illustrative example. Consider the principle \(\lem*\),
  which says \(\for{p:Ω}{p ∨ ¬p}\). Analysing this we get
  \begin{align*}
    \i{\for{p:Ω}{p ∨ ¬p}}
    = ⋀_{p∈\h} {\e p ⇒ \i{p ∨ ¬p}}\\
    = ⋀_{p∈\h} {\p{p ⇔ p} ⇒ p ∨ \p{p ⇒ ⊥}}\\
    = ⋀_{p∈\h} {⊤ ⇒ p ∨ p^⊥}\\
    = ⋀_{p∈\h} {p ∨ p^⊥}
  \end{align*}

  The principle will hold in \(\h\) precisely when this interpretation equals
  \(⊤\). Since it is a conjunction that will hold only when all of its conjuncts
  are \(⊤\). But that will only happen if every element of \(p ∈ \h\) satisfies
  \(p ∨ p^⊥ = ⊤\), or in other words, when every element of the Heyting algebra
  is complemented.
\end{example}

\begin{theorem}[Folklore]
  The \(\lem*\) holds over \(\h\) iff \(\h\) is a Boolean algebra.
\end{theorem}

We can show a weaker version the same way.

\begin{exercise}
  Show that \(\wlem*\) holds over \(\h\) iff in \(\h\) every pseudo-complement is
  complemented.

  In the case where \(\h\) is a topology this corresponds to extremally
  disconnected spaces (closure of every open set is open).
\end{exercise}


\subsubsection{Intermezzo: Real numbers}

Let \(X\) be some topological space. Then we can synthesise it to a
Heyting-valued set.

\begin{construction}
  The set \(\set{f : 𝒞(p, X)}{p ∈ \h}\) with equality
  \[ \eq f g ≔ ⋁\set{q ≤ \e f ∧ \e g}{f\res q = g\res q} \]
  is the Heyting-valued set \(X_{\h}\).
\end{construction}

We now have two synthesisations of \(ℝ\). One is the discrete one \(\c ℝ\), and
the other the topological one \(ℝ_{\h}\). It turns out the latter one is the
object of Dedekind reals.

TODO: proof

\begin{example}
  The analytic limited principle of omniscience (\(\alpo*\)) states that the apartness on
  real numbers is decidable, or \(\for{x : ℝ}{x = 0 ∨ x \apart 0}\).

  Apartness is defined as \(x \apart y ≔ x < y ∨ x > y\) and
  \(x > y ≔ \exist{q:ℚ}{x > q > y}\), where comparisons with rationals are done
  by containment in the cuts.
  Then \(\alpo*\) computes to ``for every continuous function \(f : 𝒞(p, ℝ)\)
  \({⊤ = \i{f = 0 ∨ f \apart 0}}\)''. This further computes to
  \({⊤ = \p{f^{-1}(ℝ⧵\{0\})}^⊥ ∨ f^{-1}(ℝ⧵\{0\})}\).
  Then \(\alpo*\) is satisfied if and only if for every continuous function
  \(f : 𝒞(p, ℝ)\) \(f^{-1}(ℝ⧵\{0\})\) is complemented.

  In the topological case, we can rephrase this to say for every continuous map
  \(f : 𝒞(U, ℝ)\) the zero-set \(f^{-1}\{0\}\) is open.
\end{example}

A similar computation can be carried out for the Cantor space \({2^ℕ}_{\h}\). It
turns out this space is just the object \(2^ℕ\).

\begin{exercise}
  \sloppy
  Let \(\h\) be a second-countable topological space (or a localic equivalent).
  Show that \(\lpo*\) holds over \(\h\) iff \(\h\) is locally connected.

  Note: This will need some topological computation, but a useful lemma is that
  a space is locally connected iff intersection of relatively clopen subsets of
  \(U\) is again relatively clopen in \(U\). You may take this as given (it is
  not part of the exercise per se, but it is what needs to be done in practice
  to obtain useful characterisations).
\end{exercise}



%%% Local Variables:
%%% mode: latex
%%% TeX-master: "../main"
%%% End:
